\section{Introduction}
\label{sec:intro}

% OUTLINE (prose to be written with /paper-language, section-by-section):
%
% P1 -- The problem, in the challenge's own terms. An engineer receives a
%   natural-language issue report and logs. No failing test, no debugger,
%   no reproducible execution. They must find the defective code region.
%   Ground with one real benchmark example (a MariaDB or Kafka issue where
%   the stack trace names the WRONG file).
%
% P2 -- Why existing routes fall short (mirrors the challenge text):
%   test-based fault localization needs execution evidence that does not
%   exist at report time; AIOps RCA ranks services/components, not code
%   regions; codebase-retrieval methods treat every issue as the first
%   issue ever filed.
%
% P3 -- The idea. Software failures recur. The incidents a team already
%   resolved left two artifacts behind: a report describing symptoms, and
%   a fix commit naming the code that was actually at fault. We treat the
%   resolved history as INCIDENT MEMORY: past incidents similar to the new
%   report vote for the files their fixes touched. This signal is fused
%   with log-to-code retrieval.
%
% P4 -- The measurement problem and our benchmark: no public dataset pairs
%   issue + logs + code snapshot + ground-truth faulty regions. We mine one
%   from two public datasets x git history: \nInstances real incidents
%   across \nProjects organizations (Kafka, MariaDB, MongoDB, Qt), gold =
%   files modified by the fix commit, snapshot = parent of the fix.
%   Stratified by whether the report names a gold file (\nHardPct hard).
%
% P5 -- Results preview + contributions list:
%   C1 benchmark (released); C2 memory-augmented localizer, beats the
%   strongest baseline on all four projects (acc@5 +0.057..+0.175, all
%   significant after BH); C3 the memory finding: on the hard stratum
%   memory lifts fusion on ALL FOUR projects (+0.080..+0.242 acc@5, all
%   SIG), and its value scales with the density of linked history;
%   C4 honest negatives (where fusion still loses to content retrieval).
